% ============================================================
\chapter{Iterative Methods for Linear Systems}
\label{ch:iterative}
% ============================================================

\section{Motivating example}
In finite element methods, the matrix $A$ is sparse ($\BigO(n)$ nonzero entries out of $n^2$). LU costs $\BigO(n^3)$, which is unacceptable for $n=10^6$. Iterative methods exploit the sparse structure: each iteration costs $\BigO(n)$.

\section{General principle}

\begin{definition}[Splitting method]
We write $A=M-N$ with $M$ easily invertible. The iteration is:
\[
Mx^{(k+1)} = Nx^{(k)}+b, \quad\text{i.e.}\quad x^{(k+1)}=Gx^{(k)}+c,
\]
where $G=M^{-1}N$ is the \textbf{iteration matrix} and $c=M^{-1}b$.
\end{definition}

\begin{theorem}[Convergence]\label{thm:iter-conv}
The method converges for every $x^{(0)}$ if and only if $\rho(G)<1$ (spectral radius).
\end{theorem}

\begin{proof}
$e^{(k)}=x^{(k)}-x^*=G^k e^{(0)}$. Now $\|G^k\|\to 0 \iff \rho(G)<1$.
\end{proof}

\section{Jacobi method}\index{Jacobi}

\begin{definition}
$A=D-E-F$ (diagonal, lower triangular, upper triangular). $M=D$:
\[
x_i^{(k+1)} = \frac{1}{a_{ii}}\left(b_i-\sum_{j\neq i}a_{ij}x_j^{(k)}\right).
\]
Iteration matrix: $G_J=D^{-1}(E+F)$.
\end{definition}

\section{Gauss--Seidel method}\index{Gauss--Seidel}

\begin{definition}
$M=D-E$:
\[
x_i^{(k+1)} = \frac{1}{a_{ii}}\left(b_i-\sum_{j<i}a_{ij}x_j^{(k+1)}-\sum_{j>i}a_{ij}x_j^{(k)}\right).
\]
Iteration matrix: $G_{GS}=(D-E)^{-1}F$.
\end{definition}

\begin{theorem}[Convergence for SPD matrices]
If $A$ is SPD, Gauss--Seidel converges.
\end{theorem}

\begin{theorem}[Convergence for diagonally dominant matrices]
If $A$ is strictly diagonally dominant ($|a_{ii}|>\sum_{j\neq i}|a_{ij}|$), then both Jacobi and Gauss--Seidel converge.
\end{theorem}

\section{Successive over-relaxation (SOR)}\index{SOR}

\begin{definition}
We introduce a parameter $\omega\in(0,2)$:
\[
x_i^{(k+1)} = (1-\omega)x_i^{(k)} + \frac{\omega}{a_{ii}}\left(b_i-\sum_{j<i}a_{ij}x_j^{(k+1)}-\sum_{j>i}a_{ij}x_j^{(k)}\right).
\]
$\omega=1$: Gauss--Seidel. $\omega>1$: over-relaxation. $\omega<1$: under-relaxation.
\end{definition}

\begin{theorem}
SOR converges if and only if $0<\omega<2$. The optimal $\omega$ depends on $\rho(G_J)$.
\end{theorem}

\section{Conjugate gradient method}\index{conjugate gradient}

\begin{definition}
For $A$ SPD, solving $Ax=b$ is equivalent to minimising $\phi(x)=\frac{1}{2}x^\top Ax-b^\top x$.
\end{definition}

\begin{algorithm}[H]
\caption{Conjugate gradient}\label{alg:cg}
\KwInput{$A$ SPD, $b$, $x_0$, tolerance $\eps$}
\KwOutput{$x\approx A^{-1}b$}
$r_0\gets b-Ax_0$, $p_0\gets r_0$\;
\For{$k=0,1,2,\ldots$}{
  $\alpha_k\gets\frac{r_k^\top r_k}{p_k^\top Ap_k}$\;
  $x_{k+1}\gets x_k+\alpha_k p_k$\;
  $r_{k+1}\gets r_k-\alpha_k Ap_k$\;
  \If{$\|r_{k+1}\|<\eps$}{\Return{$x_{k+1}$}}
  $\beta_k\gets\frac{r_{k+1}^\top r_{k+1}}{r_k^\top r_k}$\;
  $p_{k+1}\gets r_{k+1}+\beta_k p_k$\;
}
\end{algorithm}

\begin{theorem}[Convergence of the conjugate gradient method]
In exact arithmetic, CG converges in at most $n$ iterations. The error satisfies:
\[
\|e_k\|_A \leq 2\left(\frac{\sqrt{\kappa}-1}{\sqrt{\kappa}+1}\right)^k\|e_0\|_A, \quad \kappa=\cond_2(A).
\]
\end{theorem}

\section{Implementation}

\begin{codeblock}[title=\textbf{Python}]
\begin{minted}{python}
import numpy as np

def jacobi(A, b, x0, tol=1e-10, maxiter=1000):
    D = np.diag(np.diag(A))
    R = A - D
    x = x0.copy()
    for k in range(maxiter):
        x_new = np.linalg.solve(D, b - R @ x)
        if np.linalg.norm(x_new - x) < tol:
            return x_new, k+1
        x = x_new
    return x, maxiter

def conjugate_gradient(A, b, x0, tol=1e-10, maxiter=1000):
    x = x0.copy()
    r = b - A @ x
    p = r.copy()
    rs_old = r @ r
    for k in range(maxiter):
        Ap = A @ p
        alpha = rs_old / (p @ Ap)
        x += alpha * p
        r -= alpha * Ap
        rs_new = r @ r
        if np.sqrt(rs_new) < tol:
            return x, k+1
        p = r + (rs_new / rs_old) * p
        rs_old = rs_new
    return x, maxiter

# Test : matrice tridiagonale
n = 100
A = np.diag(2*np.ones(n)) - np.diag(np.ones(n-1), 1) - np.diag(np.ones(n-1), -1)
b = np.ones(n)
x0 = np.zeros(n)

x_j, nj = jacobi(A, b, x0)
x_cg, ncg = conjugate_gradient(A, b, x0)
print(f"Jacobi: {nj} iterations")
print(f"CG:     {ncg} iterations")
\end{minted}
\end{codeblock}

\begin{codeblock}[title=\textbf{Julia}]
\begin{minted}{julia}
function cg(A, b, x0; tol=1e-10, maxiter=1000)
    x = copy(x0)
    r = b - A * x
    p = copy(r)
    rs = dot(r, r)
    for k in 1:maxiter
        Ap = A * p
        alpha = rs / dot(p, Ap)
        x .+= alpha .* p
        r .-= alpha .* Ap
        rs_new = dot(r, r)
        sqrt(rs_new) < tol && return x, k
        p .= r .+ (rs_new / rs) .* p
        rs = rs_new
    end
    return x, maxiter
end
\end{minted}
\end{codeblock}

\section{Exercises}

\begin{exercise}[$\star$]
Apply 3 iterations of Jacobi and Gauss--Seidel to the system $\begin{pmatrix}4&1\\1&3\end{pmatrix}x=\begin{pmatrix}1\\2\end{pmatrix}$, $x^{(0)}=(0,0)^\top$.
\end{exercise}

\begin{exercise}[$\star\star$]
Show that for a tridiagonal matrix $T_n$ of size $n$, $\rho(G_J)=\cos(\pi/(n+1))$. Deduce the number of Jacobi iterations needed for a precision $\eps$.
\end{exercise}

\begin{exercise}[$\star\star\star$ -- Project]
Compare Jacobi, Gauss--Seidel, SOR and CG for the 2D discretisation of the Laplacian $-\Delta u=f$ on $[0,1]^2$. Plot the number of iterations as a function of $n$.
\end{exercise}

\begin{keyformulas}
\begin{align*}
\text{Jacobi} &: x^{(k+1)}=D^{-1}(b-(A-D)x^{(k)}) \\
\text{Gauss--Seidel} &: G=(D-E)^{-1}F \\
\text{CG} &: \|e_k\|_A\leq 2\left(\frac{\sqrt{\kappa}-1}{\sqrt{\kappa}+1}\right)^k\|e_0\|_A
\end{align*}
\end{keyformulas}
